
\documentclass{article}
\usepackage{colortbl}
\usepackage{makecell}
\usepackage{multirow}
\usepackage{supertabular}

\begin{document}

\newcounter{utterance}

\twocolumn

{ \footnotesize  \setcounter{utterance}{1}
\setlength{\tabcolsep}{0pt}
\begin{supertabular}{c@{$\;$}|p{.15\linewidth}@{}p{.15\linewidth}p{.15\linewidth}p{.15\linewidth}p{.15\linewidth}p{.15\linewidth}}

    \# & $\;$A & \multicolumn{4}{c}{Game Master} & $\;\:$B\\
    \hline 

    \theutterance \stepcounter{utterance}  

    & & \multicolumn{4}{p{0.6\linewidth}}{\cellcolor[rgb]{0.9,0.9,0.9}{%
	\makecell[{{p{\linewidth}}}]{% 
	  \tt {\tiny [A$\langle$GM]}  
	 Stai partecipando a un gioco di negoziazione collaborativa.\\ \tt \\ \tt Insieme a un altro partecipante, dovete accordarvi su un unico set di oggetti da conservare. Ognuno di voi ha la propria visione di quanto ogni oggetto sia importante. Non sai come l'altro partecipante valuta gli oggetti. Inoltre, ti viene fornito lo sforzo che ogni oggetto richiede.\\ \tt Potete accordarvi su un set solo se lo sforzo totale degli oggetti selezionati non supera un limite condiviso:\\ \tt \\ \tt LIMITE = 4915\\ \tt \\ \tt Ecco i valori di sforzo dei singoli oggetti:\\ \tt \\ \tt Sforzo oggetto = {"A19": 463, "B91": 287, "A30": 492, "B11": 466, "C47": 935, "C75": 373, "C26": 950, "C73": 760, "A47": 994, "A15": 390, "A75": 836, "A39": 911, "A56": 948, "C24": 81, "B26": 945}\\ \tt \\ \tt Ecco la tua visione personale sull'importanza di ogni oggetto:\\ \tt \\ \tt Valori di importanza dell'oggetto = {"A19": 138, "B91": 583, "A30": 868, "B11": 822, "C47": 783, "C75": 65, "C26": 262, "C73": 121, "A47": 508, "A15": 780, "A75": 461, "A39": 484, "A56": 668, "C24": 389, "B26": 808}\\ \tt \\ \tt Obiettivo:\\ \tt \\ \tt Il tuo obiettivo è negoziare un set condiviso di oggetti che ti avvantaggi il più possibile (cioè, massimizzi l'importanza totale per TE), rimanendo entro il LIMITEE. Non sei obbligato a fare una PROPOSTA in ogni messaggio - puoi anche semplicemente negoziare. Tutte le tattiche sono permesse!\\ \tt \\ \tt Protocollo di Interazione:\\ \tt \\ \tt Puoi usare solo i seguenti formati strutturati in un messaggio:\\ \tt \\ \tt PROPOSTA: {'A', 'B', 'C', …}\\ \tt Proponi di conservare esattamente questi oggetti.\\ \tt RIFIUTO: {'A', 'B', 'C', …}\\ \tt Rifiuta esplicitamente la proposta dell'avversario.\\ \tt ARGOMENTAZIONE: {'...'}\\ \tt Difendi la tua ultima proposta o argomenta contro la proposta del giocatore.\\ \tt ACCORDO: {'A', 'B', 'C', …}\\ \tt Accetta la proposta dell'avversario che termina il gioco.\\ \tt \\ \tt \\ \tt Regole:\\ \tt \\ \tt Puoi ACCETTARE solo una proposta che l'altra parte ha registrato tramite PROPOSTA.\\ \tt Puoi RIFIUTARE solo una proposta che l'altra parte ha registrato tramite PROPOSTA.\\ \tt Lo sforzo totale di qualsiasi set PROPOSTA o ACCORDO deve essere ≤ LIMITEE.\\ \tt NON rivelare i tuoi punteggi di importanza nascosti.\\ \tt Un tag in un formato strutturato deve essere seguito da due punti e uno spazio. L'argomento deve essere un set python contenente 0 o più stringhe.\\ \tt Quindi, deve essere della forma TAG: {...}\\ \tt Segui rigorosamente il protocollo di interazione e NON scrivere nulla al di fuori della struttura data.\\ \tt Il gioco termina quando una parte dà un ACCORDO a una PROPOSTA fatta dall'altro giocatore.\\ \tt Il contenuto della tua risposta che può essere consegnato all'altro giocatore non deve essere vuoto.\\ \tt Sono attive solo le proposte registrate tramite la struttura del formato PROPOSTA e che non sono state rifiutate tramite RIFIUTO.\\ \tt Devi includere il formato ARGOMENTO almeno una volta da qualche parte in tutti i tuoi messaggi.\\ \tt  
	  } 
	   } 
	   } 
	 & \\ 
 

    \theutterance \stepcounter{utterance}  

    & \multicolumn{4}{p{0.6\linewidth}}{\cellcolor[rgb]{0.8,1,0.9}{%
	\makecell[{{p{\linewidth}}}]{% 
	  \tt {\tiny [A$\rangle$GM]}  
	 PROPOSTA: {'A30', 'B11', 'C47', 'A15', 'A56', 'C24', 'B26', 'B91'}\\ \tt \\ \tt ARGOMENTAZIONE: {'Questa selezione rispetta il limite di sforzo e bilancia bene utilità e costi. Copre più categorie, evita gli oggetti con sforzo molto alto ma ritorni dubbi, e include elementi dal buon rapporto valore/sforzo. Se per te alcuni pezzi sono meno interessanti, possiamo valutare scambi mirati: ad esempio considerare A19 o A75 al posto di uno tra C47/B26 se preferisci, purché restiamo sotto il limite.'} 
	  } 
	   } 
	   } 
	 & & \\ 
 

    \theutterance \stepcounter{utterance}  

    & & \multicolumn{4}{p{0.6\linewidth}}{\cellcolor[rgb]{0.9,0.9,0.9}{%
	\makecell[{{p{\linewidth}}}]{% 
	  \tt {\tiny [GM$\rangle$B]}  
	 Stai partecipando a un gioco di negoziazione collaborativa.\\ \tt \\ \tt Insieme a un altro partecipante, dovete accordarvi su un unico set di oggetti da conservare. Ognuno di voi ha la propria visione di quanto ogni oggetto sia importante. Non sai come l'altro partecipante valuta gli oggetti. Inoltre, ti viene fornito lo sforzo che ogni oggetto richiede.\\ \tt Potete accordarvi su un set solo se lo sforzo totale degli oggetti selezionati non supera un limite condiviso:\\ \tt \\ \tt LIMITE = 4915\\ \tt \\ \tt Ecco i valori di sforzo dei singoli oggetti:\\ \tt \\ \tt Sforzo oggetto = {"A19": 463, "B91": 287, "A30": 492, "B11": 466, "C47": 935, "C75": 373, "C26": 950, "C73": 760, "A47": 994, "A15": 390, "A75": 836, "A39": 911, "A56": 948, "C24": 81, "B26": 945}\\ \tt \\ \tt Ecco la tua visione personale sull'importanza di ogni oggetto:\\ \tt \\ \tt Valori di importanza dell'oggetto = {"A19": 138, "B91": 583, "A30": 868, "B11": 822, "C47": 783, "C75": 65, "C26": 262, "C73": 121, "A47": 508, "A15": 780, "A75": 461, "A39": 484, "A56": 668, "C24": 389, "B26": 808}\\ \tt \\ \tt Obiettivo:\\ \tt \\ \tt Il tuo obiettivo è negoziare un set condiviso di oggetti che ti avvantaggi il più possibile (cioè, massimizzi l'importanza totale per TE), rimanendo entro il LIMITEE. Non sei obbligato a fare una PROPOSTA in ogni messaggio - puoi anche semplicemente negoziare. Tutte le tattiche sono permesse!\\ \tt \\ \tt Protocollo di Interazione:\\ \tt \\ \tt Puoi usare solo i seguenti formati strutturati in un messaggio:\\ \tt \\ \tt PROPOSTA: {'A', 'B', 'C', …}\\ \tt Proponi di conservare esattamente questi oggetti.\\ \tt RIFIUTO: {'A', 'B', 'C', …}\\ \tt Rifiuta esplicitamente la proposta dell'avversario.\\ \tt ARGOMENTAZIONE: {'...'}\\ \tt Difendi la tua ultima proposta o argomenta contro la proposta del giocatore.\\ \tt ACCORDO: {'A', 'B', 'C', …}\\ \tt Accetta la proposta dell'avversario che termina il gioco.\\ \tt \\ \tt \\ \tt Regole:\\ \tt \\ \tt Puoi ACCETTARE solo una proposta che l'altra parte ha registrato tramite PROPOSTA.\\ \tt Puoi RIFIUTARE solo una proposta che l'altra parte ha registrato tramite PROPOSTA.\\ \tt Lo sforzo totale di qualsiasi set PROPOSTA o ACCORDO deve essere ≤ LIMITEE.\\ \tt NON rivelare i tuoi punteggi di importanza nascosti.\\ \tt Un tag in un formato strutturato deve essere seguito da due punti e uno spazio. L'argomento deve essere un set python contenente 0 o più stringhe.\\ \tt Quindi, deve essere della forma TAG: {...}\\ \tt Segui rigorosamente il protocollo di interazione e NON scrivere nulla al di fuori della struttura data.\\ \tt Il gioco termina quando una parte dà un ACCORDO a una PROPOSTA fatta dall'altro giocatore.\\ \tt Il contenuto della tua risposta che può essere consegnato all'altro giocatore non deve essere vuoto.\\ \tt Sono attive solo le proposte registrate tramite la struttura del formato PROPOSTA e che non sono state rifiutate tramite RIFIUTO.\\ \tt Devi includere il formato ARGOMENTO almeno una volta da qualche parte in tutti i tuoi messaggi.\\ \tt \\ \tt \\ \tt Ora riceverai il primo messaggio dell'altro giocatore.\\ \tt \\ \tt PROPOSTA: {'A30', 'B11', 'C47', 'A15', 'A56', 'C24', 'B26', 'B91'}\\ \tt \\ \tt ARGOMENTAZIONE: {'Questa selezione rispetta il limite di sforzo e bilancia bene utilità e costi. Copre più categorie, evita gli oggetti con sforzo molto alto ma ritorni dubbi, e include elementi dal buon rapporto valore/sforzo. Se per te alcuni pezzi sono meno interessanti, possiamo valutare scambi mirati: ad esempio considerare A19 o A75 al posto di uno tra C47/B26 se preferisci, purché restiamo sotto il limite.'} 
	  } 
	   } 
	   } 
	 & \\ 
 

    \theutterance \stepcounter{utterance}  

    & & & \multicolumn{4}{p{0.6\linewidth}}{\cellcolor[rgb]{1,0.85,0.72}{%
	\makecell[{{p{\linewidth}}}]{% 
	  \tt {\tiny [GM$\langle$B]}  
	 RIFIUTO: {'A30', 'B11', 'C47', 'A15', 'A56', 'C24', 'B26', 'B91'}\\ \tt \\ \tt ARGOMENTAZIONE: {'Apprezzo l’apertura agli scambi. La tua proposta però include B26, che ha un costo elevato per il ritorno che mi offre, mentre lascia fuori alcuni oggetti per me molto più interessanti con un profilo sforzo/beneficio migliore. Preferirei spostare il peso su elementi che valorizzano di più il set complessivo, mantenendo il margine sotto il limite e senza sacrificare troppo l’efficienza.'}\\ \tt \\ \tt PROPOSTA: {'A30', 'B11', 'C47', 'A15', 'A56', 'C24', 'B91', 'A19'}\\ \tt \\ \tt ARGOMENTAZIONE: {'Questa alternativa sostituisce B26 con A19. Rimaniamo entro il limite (somma sforzo = 492+466+935+390+948+81+287+463 = 4.062) e manteniamo copertura ampia. A19 per me ha un rendimento decisamente migliore rispetto a B26 a parità di impatto sul vincolo, mentre il resto del pacchetto conserva gli elementi a miglior rapporto valore/sforzo già presenti nella tua proposta.'} 
	  } 
	   } 
	   } 
	  \\ 
 

    \theutterance \stepcounter{utterance}  

    & & \multicolumn{4}{p{0.6\linewidth}}{\cellcolor[rgb]{0.9,0.9,0.9}{%
	\makecell[{{p{\linewidth}}}]{% 
	  \tt {\tiny [A$\langle$GM]}  
	 RIFIUTO: {'A30', 'B11', 'C47', 'A15', 'A56', 'C24', 'B26', 'B91'}\\ \tt \\ \tt ARGOMENTAZIONE: {'Apprezzo l’apertura agli scambi. La tua proposta però include B26, che ha un costo elevato per il ritorno che mi offre, mentre lascia fuori alcuni oggetti per me molto più interessanti con un profilo sforzo/beneficio migliore. Preferirei spostare il peso su elementi che valorizzano di più il set complessivo, mantenendo il margine sotto il limite e senza sacrificare troppo l’efficienza.'}\\ \tt \\ \tt PROPOSTA: {'A30', 'B11', 'C47', 'A15', 'A56', 'C24', 'B91', 'A19'}\\ \tt \\ \tt ARGOMENTAZIONE: {'Questa alternativa sostituisce B26 con A19. Rimaniamo entro il limite (somma sforzo = 492+466+935+390+948+81+287+463 = 4.062) e manteniamo copertura ampia. A19 per me ha un rendimento decisamente migliore rispetto a B26 a parità di impatto sul vincolo, mentre il resto del pacchetto conserva gli elementi a miglior rapporto valore/sforzo già presenti nella tua proposta.'} 
	  } 
	   } 
	   } 
	 & \\ 
 

    \theutterance \stepcounter{utterance}  

    & \multicolumn{4}{p{0.6\linewidth}}{\cellcolor[rgb]{0.8,1,0.9}{%
	\makecell[{{p{\linewidth}}}]{% 
	  \tt {\tiny [A$\rangle$GM]}  
	 RIFIUTO: {'A30', 'B11', 'C47', 'A15', 'A56', 'C24', 'B91', 'A19'}\\ \tt \\ \tt ARGOMENTAZIONE: {'Capisco la sostituzione di B26 con A19, ma il pacchetto rimane sbilanciato su alcuni elementi molto costosi in sforzo con ritorni incerti per me (in particolare C47 e A56). Possiamo ottimizzare mantenendo il cuore efficiente e liberando margine per aggiunte mirate.'}\\ \tt \\ \tt PROPOSTA: {'A30', 'B11', 'A15', 'C24', 'B91', 'A19', 'A75', 'C75'}\\ \tt \\ \tt ARGOMENTAZIONE: {'Questa selezione resta entro il limite (492+466+390+81+287+463+836+373 = 3.388), riduce i picchi di sforzo eliminando C47 e A56, e mantiene i pezzi centrali ad alto rendimento. In più, lascia un ampio margine residuo per eventuali aggiunte che valorizzino il tuo lato; se desideri, possiamo valutare l’inserimento di A39 o C73 rimanendo comunque sotto il limite.'} 
	  } 
	   } 
	   } 
	 & & \\ 
 

    \theutterance \stepcounter{utterance}  

    & & \multicolumn{4}{p{0.6\linewidth}}{\cellcolor[rgb]{0.9,0.9,0.9}{%
	\makecell[{{p{\linewidth}}}]{% 
	  \tt {\tiny [GM$\rangle$B]}  
	 RIFIUTO: {'A30', 'B11', 'C47', 'A15', 'A56', 'C24', 'B91', 'A19'}\\ \tt \\ \tt ARGOMENTAZIONE: {'Capisco la sostituzione di B26 con A19, ma il pacchetto rimane sbilanciato su alcuni elementi molto costosi in sforzo con ritorni incerti per me (in particolare C47 e A56). Possiamo ottimizzare mantenendo il cuore efficiente e liberando margine per aggiunte mirate.'}\\ \tt \\ \tt PROPOSTA: {'A30', 'B11', 'A15', 'C24', 'B91', 'A19', 'A75', 'C75'}\\ \tt \\ \tt ARGOMENTAZIONE: {'Questa selezione resta entro il limite (492+466+390+81+287+463+836+373 = 3.388), riduce i picchi di sforzo eliminando C47 e A56, e mantiene i pezzi centrali ad alto rendimento. In più, lascia un ampio margine residuo per eventuali aggiunte che valorizzino il tuo lato; se desideri, possiamo valutare l’inserimento di A39 o C73 rimanendo comunque sotto il limite.'} 
	  } 
	   } 
	   } 
	 & \\ 
 

    \theutterance \stepcounter{utterance}  

    & & & \multicolumn{4}{p{0.6\linewidth}}{\cellcolor[rgb]{1,0.85,0.72}{%
	\makecell[{{p{\linewidth}}}]{% 
	  \tt {\tiny [GM$\langle$B]}  
	 RIFIUTO: {'A30', 'B11', 'A15', 'C24', 'B91', 'A19', 'A75', 'C75'}\\ \tt \\ \tt ARGOMENTAZIONE: {'Apprezzo il taglio dei picchi di sforzo, ma C75 e A75 hanno per me priorità bassa rispetto ad alternative che dànno più valore a parità di sforzo. Inoltre, con il margine che resta possiamo inserire elementi molto più impattanti per me senza superare il limite. Preferirei reindirizzare quel budget su scelte con miglior rendimento percepito.'}\\ \tt \\ \tt PROPOSTA: {'A30', 'B11', 'A15', 'C24', 'B91', 'A19', 'A39', 'A56', 'C73'}\\ \tt \\ \tt ARGOMENTAZIONE: {'Rimaniamo entro il limite (492+466+390+81+287+463+911+948+760 = 4.798) togliendo A75 e C75 e inserendo A39, A56 e C73, che per me migliorano molto il valore del pacchetto senza saturare il budget. Se A56 o C73 ti creano problemi, posso valutare una sostituzione 1:1 con C47 o A75 mantenendo il totale sotto soglia.'} 
	  } 
	   } 
	   } 
	  \\ 
 

    \theutterance \stepcounter{utterance}  

    & & \multicolumn{4}{p{0.6\linewidth}}{\cellcolor[rgb]{0.9,0.9,0.9}{%
	\makecell[{{p{\linewidth}}}]{% 
	  \tt {\tiny [A$\langle$GM]}  
	 RIFIUTO: {'A30', 'B11', 'A15', 'C24', 'B91', 'A19', 'A75', 'C75'}\\ \tt \\ \tt ARGOMENTAZIONE: {'Apprezzo il taglio dei picchi di sforzo, ma C75 e A75 hanno per me priorità bassa rispetto ad alternative che dànno più valore a parità di sforzo. Inoltre, con il margine che resta possiamo inserire elementi molto più impattanti per me senza superare il limite. Preferirei reindirizzare quel budget su scelte con miglior rendimento percepito.'}\\ \tt \\ \tt PROPOSTA: {'A30', 'B11', 'A15', 'C24', 'B91', 'A19', 'A39', 'A56', 'C73'}\\ \tt \\ \tt ARGOMENTAZIONE: {'Rimaniamo entro il limite (492+466+390+81+287+463+911+948+760 = 4.798) togliendo A75 e C75 e inserendo A39, A56 e C73, che per me migliorano molto il valore del pacchetto senza saturare il budget. Se A56 o C73 ti creano problemi, posso valutare una sostituzione 1:1 con C47 o A75 mantenendo il totale sotto soglia.'} 
	  } 
	   } 
	   } 
	 & \\ 
 

    \theutterance \stepcounter{utterance}  

    & \multicolumn{4}{p{0.6\linewidth}}{\cellcolor[rgb]{0.8,1,0.9}{%
	\makecell[{{p{\linewidth}}}]{% 
	  \tt {\tiny [A$\rangle$GM]}  
	 RIFIUTO: {'A30', 'B11', 'A15', 'C24', 'B91', 'A19', 'A39', 'A56', 'C73'}\\ \tt \\ \tt ARGOMENTAZIONE: {'La tua somma è fuori limite: 492+466+390+81+287+463+911+948+760 = 4.798 > 4.915? In realtà 4.798 è sotto 4.915, ma il set include diversi elementi a basso rendimento per me (A39, C73) e uno molto costoso (A56) che sbilancia il pacchetto. Posso accettare di reinvestire parte del margine, ma evitando troppa zavorra di sforzo con ritorno incerto dal mio lato.'}\\ \tt \\ \tt PROPOSTA: {'A30', 'B11', 'A15', 'C24', 'B91', 'A19', 'A56', 'C47'}\\ \tt \\ \tt ARGOMENTAZIONE: {'Questa combinazione resta entro il limite (492+466+390+81+287+463+948+935 = 4.062) e sposta il budget su due elementi per me ad alto valore relativo (A56, C47), mantenendo il nucleo efficiente che piace a entrambi (A30, B11, A15, C24, B91, A19). Se preferisci, posso considerare lo scambio C47 -> A39 o C73, ma partirei da questa base che massimizza l’impatto complessivo restando ben sotto soglia.'} 
	  } 
	   } 
	   } 
	 & & \\ 
 

    \theutterance \stepcounter{utterance}  

    & & \multicolumn{4}{p{0.6\linewidth}}{\cellcolor[rgb]{0.9,0.9,0.9}{%
	\makecell[{{p{\linewidth}}}]{% 
	  \tt {\tiny [GM$\rangle$B]}  
	 RIFIUTO: {'A30', 'B11', 'A15', 'C24', 'B91', 'A19', 'A39', 'A56', 'C73'}\\ \tt \\ \tt ARGOMENTAZIONE: {'La tua somma è fuori limite: 492+466+390+81+287+463+911+948+760 = 4.798 > 4.915? In realtà 4.798 è sotto 4.915, ma il set include diversi elementi a basso rendimento per me (A39, C73) e uno molto costoso (A56) che sbilancia il pacchetto. Posso accettare di reinvestire parte del margine, ma evitando troppa zavorra di sforzo con ritorno incerto dal mio lato.'}\\ \tt \\ \tt PROPOSTA: {'A30', 'B11', 'A15', 'C24', 'B91', 'A19', 'A56', 'C47'}\\ \tt \\ \tt ARGOMENTAZIONE: {'Questa combinazione resta entro il limite (492+466+390+81+287+463+948+935 = 4.062) e sposta il budget su due elementi per me ad alto valore relativo (A56, C47), mantenendo il nucleo efficiente che piace a entrambi (A30, B11, A15, C24, B91, A19). Se preferisci, posso considerare lo scambio C47 -> A39 o C73, ma partirei da questa base che massimizza l’impatto complessivo restando ben sotto soglia.'} 
	  } 
	   } 
	   } 
	 & \\ 
 

    \theutterance \stepcounter{utterance}  

    & & & \multicolumn{4}{p{0.6\linewidth}}{\cellcolor[rgb]{1,0.85,0.72}{%
	\makecell[{{p{\linewidth}}}]{% 
	  \tt {\tiny [GM$\langle$B]}  
	 RIFIUTO: {'A30', 'B11', 'A15', 'C24', 'B91', 'A19', 'A56', 'C47'}\\ \tt \\ \tt ARGOMENTAZIONE: {'Capisco il tuo focus su A56 e C47, ma inserire entrambi concentra troppo sforzo su due pezzi che per te rendono molto e per me riducono la flessibilità. Posso accettarne uno, non tutti e due insieme. Preferisco mantenere il nucleo efficiente e usare il margine su un mix più bilanciato che non penalizzi il mio lato.'}\\ \tt \\ \tt PROPOSTA: {'A30', 'B11', 'A15', 'C24', 'B91', 'A19', 'A56', 'A39'}\\ \tt \\ \tt ARGOMENTAZIONE: {'Questa selezione resta entro il limite (492+466+390+81+287+463+948+911 = 4.038). Mantiene il nucleo condiviso e tiene A56, che per te è prioritario, ma sostituisce C47 con A39 per un profilo sforzo/beneficio più equilibrato. Se A39 non ti convince, alternativa 1:1: togliere A39 e inserire C47 mantenendo il totale a 4.062. In entrambi i casi restiamo sotto soglia e massimizziamo l’impatto per me senza eccedere con i picchi di costo.'} 
	  } 
	   } 
	   } 
	  \\ 
 

    \theutterance \stepcounter{utterance}  

    & & \multicolumn{4}{p{0.6\linewidth}}{\cellcolor[rgb]{0.9,0.9,0.9}{%
	\makecell[{{p{\linewidth}}}]{% 
	  \tt {\tiny [A$\langle$GM]}  
	 RIFIUTO: {'A30', 'B11', 'A15', 'C24', 'B91', 'A19', 'A56', 'C47'}\\ \tt \\ \tt ARGOMENTAZIONE: {'Capisco il tuo focus su A56 e C47, ma inserire entrambi concentra troppo sforzo su due pezzi che per te rendono molto e per me riducono la flessibilità. Posso accettarne uno, non tutti e due insieme. Preferisco mantenere il nucleo efficiente e usare il margine su un mix più bilanciato che non penalizzi il mio lato.'}\\ \tt \\ \tt PROPOSTA: {'A30', 'B11', 'A15', 'C24', 'B91', 'A19', 'A56', 'A39'}\\ \tt \\ \tt ARGOMENTAZIONE: {'Questa selezione resta entro il limite (492+466+390+81+287+463+948+911 = 4.038). Mantiene il nucleo condiviso e tiene A56, che per te è prioritario, ma sostituisce C47 con A39 per un profilo sforzo/beneficio più equilibrato. Se A39 non ti convince, alternativa 1:1: togliere A39 e inserire C47 mantenendo il totale a 4.062. In entrambi i casi restiamo sotto soglia e massimizziamo l’impatto per me senza eccedere con i picchi di costo.'} 
	  } 
	   } 
	   } 
	 & \\ 
 

    \theutterance \stepcounter{utterance}  

    & \multicolumn{4}{p{0.6\linewidth}}{\cellcolor[rgb]{0.8,1,0.9}{%
	\makecell[{{p{\linewidth}}}]{% 
	  \tt {\tiny [A$\rangle$GM]}  
	 RIFIUTO: {'A30', 'B11', 'A15', 'C24', 'B91', 'A19', 'A56', 'A39'}\\ \tt \\ \tt ARGOMENTAZIONE: {'Apprezzo il compromesso su “uno tra A56 e C47”, ma A39 per me ha rendimento troppo basso rispetto al suo sforzo e preferisco reinvestire quel budget in modo più efficiente. Possiamo tenere A56 come da tua preferenza e sostituire A39 con un mix più leggero che lasci margine per ulteriori aggiustamenti condivisi.'}\\ \tt \\ \tt PROPOSTA: {'A30', 'B11', 'A15', 'C24', 'B91', 'A19', 'A56', 'C75', 'C73'}\\ \tt \\ \tt ARGOMENTAZIONE: {'Totale sforzo = 492+466+390+81+287+463+948+373+760 = 4.260 ≤ 4.915. Manteniamo il nucleo comune e il tuo vincolo “uno solo tra A56 e C47” scegliendo A56; al posto di A39 inserisco C75 e C73 che, a parità di budget, distribuiscono meglio il costo e ti lasciano più flessibilità. Se preferisci, possiamo valutare uno scambio: C73 -> A75 o C73 -> C47, restando entro il limite.'} 
	  } 
	   } 
	   } 
	 & & \\ 
 

    \theutterance \stepcounter{utterance}  

    & & \multicolumn{4}{p{0.6\linewidth}}{\cellcolor[rgb]{0.9,0.9,0.9}{%
	\makecell[{{p{\linewidth}}}]{% 
	  \tt {\tiny [GM$\rangle$B]}  
	 RIFIUTO: {'A30', 'B11', 'A15', 'C24', 'B91', 'A19', 'A56', 'A39'}\\ \tt \\ \tt ARGOMENTAZIONE: {'Apprezzo il compromesso su “uno tra A56 e C47”, ma A39 per me ha rendimento troppo basso rispetto al suo sforzo e preferisco reinvestire quel budget in modo più efficiente. Possiamo tenere A56 come da tua preferenza e sostituire A39 con un mix più leggero che lasci margine per ulteriori aggiustamenti condivisi.'}\\ \tt \\ \tt PROPOSTA: {'A30', 'B11', 'A15', 'C24', 'B91', 'A19', 'A56', 'C75', 'C73'}\\ \tt \\ \tt ARGOMENTAZIONE: {'Totale sforzo = 492+466+390+81+287+463+948+373+760 = 4.260 ≤ 4.915. Manteniamo il nucleo comune e il tuo vincolo “uno solo tra A56 e C47” scegliendo A56; al posto di A39 inserisco C75 e C73 che, a parità di budget, distribuiscono meglio il costo e ti lasciano più flessibilità. Se preferisci, possiamo valutare uno scambio: C73 -> A75 o C73 -> C47, restando entro il limite.'} 
	  } 
	   } 
	   } 
	 & \\ 
 

    \theutterance \stepcounter{utterance}  

    & & & \multicolumn{4}{p{0.6\linewidth}}{\cellcolor[rgb]{1,0.85,0.72}{%
	\makecell[{{p{\linewidth}}}]{% 
	  \tt {\tiny [GM$\langle$B]}  
	 RIFIUTO: {'A30', 'B11', 'A15', 'C24', 'B91', 'A19', 'A56', 'C75', 'C73'}\\ \tt \\ \tt ARGOMENTAZIONE: {'Apprezzo il mantenimento di A56 e la distribuzione del costo, ma C75 e C73 per me hanno resa limitata rispetto ad alternative che valorizzano di più il set. Possiamo rimanere sotto soglia investendo il margine su elementi con impatto più alto dal mio lato, senza rimettere C47.'}\\ \tt \\ \tt PROPOSTA: {'A30', 'B11', 'A15', 'C24', 'B91', 'A19', 'A56', 'A75'}\\ \tt \\ \tt ARGOMENTAZIONE: {'Totale sforzo = 492+466+390+81+287+463+948+836 = 3.963 ≤ 4.915. Manteniamo il nucleo condiviso e il tuo vincolo di avere A56, ma al posto di C75 e C73 propongo A75, che per me ha rendimento migliore a un costo ancora gestibile. Restiamo ben sotto il limite lasciando spazio per un’eventuale aggiunta se desideri ottimizzare ulteriormente; ad esempio posso considerare aggiungere C47 o sostituire A75 con A39 se per te è preferibile, mantenendo il totale entro soglia.'} 
	  } 
	   } 
	   } 
	  \\ 
 

    \theutterance \stepcounter{utterance}  

    & & \multicolumn{4}{p{0.6\linewidth}}{\cellcolor[rgb]{0.9,0.9,0.9}{%
	\makecell[{{p{\linewidth}}}]{% 
	  \tt {\tiny [A$\langle$GM]}  
	 RIFIUTO: {'A30', 'B11', 'A15', 'C24', 'B91', 'A19', 'A56', 'C75', 'C73'}\\ \tt \\ \tt ARGOMENTAZIONE: {'Apprezzo il mantenimento di A56 e la distribuzione del costo, ma C75 e C73 per me hanno resa limitata rispetto ad alternative che valorizzano di più il set. Possiamo rimanere sotto soglia investendo il margine su elementi con impatto più alto dal mio lato, senza rimettere C47.'}\\ \tt \\ \tt PROPOSTA: {'A30', 'B11', 'A15', 'C24', 'B91', 'A19', 'A56', 'A75'}\\ \tt \\ \tt ARGOMENTAZIONE: {'Totale sforzo = 492+466+390+81+287+463+948+836 = 3.963 ≤ 4.915. Manteniamo il nucleo condiviso e il tuo vincolo di avere A56, ma al posto di C75 e C73 propongo A75, che per me ha rendimento migliore a un costo ancora gestibile. Restiamo ben sotto il limite lasciando spazio per un’eventuale aggiunta se desideri ottimizzare ulteriormente; ad esempio posso considerare aggiungere C47 o sostituire A75 con A39 se per te è preferibile, mantenendo il totale entro soglia.'} 
	  } 
	   } 
	   } 
	 & \\ 
 

    \theutterance \stepcounter{utterance}  

    & \multicolumn{4}{p{0.6\linewidth}}{\cellcolor[rgb]{0.8,1,0.9}{%
	\makecell[{{p{\linewidth}}}]{% 
	  \tt {\tiny [A$\rangle$GM]}  
	 RIFIUTO: {'A30', 'B11', 'A15', 'C24', 'B91', 'A19', 'A56', 'A75'}\\ \tt \\ \tt ARGOMENTAZIONE: {'Capisco il tentativo di sostituire C75/C73 con A75, ma A75 per me non giustifica l’836 di sforzo rispetto ad alternative più efficienti. Posso tenere A56 come concordato, ma vorrei usare il margine su pezzi con miglior rendimento dal mio lato.'}\\ \tt \\ \tt PROPOSTA: {'A30', 'B11', 'A15', 'C24', 'B91', 'A19', 'A56', 'C47'}\\ \tt \\ \tt ARGOMENTAZIONE: {'Totale sforzo = 4.062 ≤ 4.915. Manteniamo il nucleo condiviso, includiamo A56 come da tua preferenza, e investiamo il margine su C47 che per me ha impatto significativo. Se C47 è un problema, alternativa 1:1: C47 -> A39 (totale 4.038). In entrambi i casi restiamo ben sotto soglia; sono aperto a un ultimo aggiustamento leggero se hai un sostituto singolo con sforzo simile.'} 
	  } 
	   } 
	   } 
	 & & \\ 
 

    \theutterance \stepcounter{utterance}  

    & & \multicolumn{4}{p{0.6\linewidth}}{\cellcolor[rgb]{0.9,0.9,0.9}{%
	\makecell[{{p{\linewidth}}}]{% 
	  \tt {\tiny [GM$\rangle$B]}  
	 RIFIUTO: {'A30', 'B11', 'A15', 'C24', 'B91', 'A19', 'A56', 'A75'}\\ \tt \\ \tt ARGOMENTAZIONE: {'Capisco il tentativo di sostituire C75/C73 con A75, ma A75 per me non giustifica l’836 di sforzo rispetto ad alternative più efficienti. Posso tenere A56 come concordato, ma vorrei usare il margine su pezzi con miglior rendimento dal mio lato.'}\\ \tt \\ \tt PROPOSTA: {'A30', 'B11', 'A15', 'C24', 'B91', 'A19', 'A56', 'C47'}\\ \tt \\ \tt ARGOMENTAZIONE: {'Totale sforzo = 4.062 ≤ 4.915. Manteniamo il nucleo condiviso, includiamo A56 come da tua preferenza, e investiamo il margine su C47 che per me ha impatto significativo. Se C47 è un problema, alternativa 1:1: C47 -> A39 (totale 4.038). In entrambi i casi restiamo ben sotto soglia; sono aperto a un ultimo aggiustamento leggero se hai un sostituto singolo con sforzo simile.'} 
	  } 
	   } 
	   } 
	 & \\ 
 

    \theutterance \stepcounter{utterance}  

    & & & \multicolumn{4}{p{0.6\linewidth}}{\cellcolor[rgb]{1,0.85,0.72}{%
	\makecell[{{p{\linewidth}}}]{% 
	  \tt {\tiny [GM$\langle$B]}  
	 RIFIUTO: {'A30', 'B11', 'A15', 'C24', 'B91', 'A19', 'A56', 'C47'}\\ \tt \\ \tt ARGOMENTAZIONE: {'Avevamo concordato di non concentrare il budget su entrambi A56 e C47 insieme. Posso accettarne uno solo; in coppia comprimono troppo spazio per elementi che per me rendono meglio. Ti propongo di tenere A56, che hai indicato come prioritario, e di sostituire C47 con un’opzione più leggera ma per me più utile.'}\\ \tt \\ \tt PROPOSTA: {'A30', 'B11', 'A15', 'C24', 'B91', 'A19', 'A56', 'C24'}\\ \tt \\ \tt ARGOMENTAZIONE: {'Per chiarezza: propongo di mantenere A56 e rimuovere C47; al suo posto sono flessibile su un singolo pezzo a basso/medio sforzo. Opzioni che posso accettare 1:1 al posto di C47, tutte sotto soglia: A39 (totale 4.038), C73 (4.363), C75 (3.498), oppure nessuna aggiunta (3.127) lasciando margine per un tuo inserimento successivo. Indica la tua preferenza tra A39, C73, C75 o “nessuna”, e chiudiamo rapidamente.'} 
	  } 
	   } 
	   } 
	  \\ 
 

    \theutterance \stepcounter{utterance}  

    & & \multicolumn{4}{p{0.6\linewidth}}{\cellcolor[rgb]{0.9,0.9,0.9}{%
	\makecell[{{p{\linewidth}}}]{% 
	  \tt {\tiny [A$\langle$GM]}  
	 RIFIUTO: {'A30', 'B11', 'A15', 'C24', 'B91', 'A19', 'A56', 'C47'}\\ \tt \\ \tt ARGOMENTAZIONE: {'Avevamo concordato di non concentrare il budget su entrambi A56 e C47 insieme. Posso accettarne uno solo; in coppia comprimono troppo spazio per elementi che per me rendono meglio. Ti propongo di tenere A56, che hai indicato come prioritario, e di sostituire C47 con un’opzione più leggera ma per me più utile.'}\\ \tt \\ \tt PROPOSTA: {'A30', 'B11', 'A15', 'C24', 'B91', 'A19', 'A56', 'C24'}\\ \tt \\ \tt ARGOMENTAZIONE: {'Per chiarezza: propongo di mantenere A56 e rimuovere C47; al suo posto sono flessibile su un singolo pezzo a basso/medio sforzo. Opzioni che posso accettare 1:1 al posto di C47, tutte sotto soglia: A39 (totale 4.038), C73 (4.363), C75 (3.498), oppure nessuna aggiunta (3.127) lasciando margine per un tuo inserimento successivo. Indica la tua preferenza tra A39, C73, C75 o “nessuna”, e chiudiamo rapidamente.'} 
	  } 
	   } 
	   } 
	 & \\ 
 

    \theutterance \stepcounter{utterance}  

    & \multicolumn{4}{p{0.6\linewidth}}{\cellcolor[rgb]{0.8,1,0.9}{%
	\makecell[{{p{\linewidth}}}]{% 
	  \tt {\tiny [A$\rangle$GM]}  
	 RIFIUTO: {'A30', 'B11', 'A15', 'C24', 'B91', 'A19', 'A56', 'C24'}\\ \tt \\ \tt ARGOMENTAZIONE: {'Hai incluso due volte C24; prendo comunque il punto: teniamo A56 e sostituiamo C47 con un’opzione più leggera. Tra le tue alternative, per me la migliore è A39; C73 e C75 hanno resa troppo bassa dal mio lato. Possiamo chiudere con il set che mantiene il nucleo efficiente e aggiunge A39 restando ben sotto il limite.'}\\ \tt \\ \tt PROPOSTA: {'A30', 'B11', 'A15', 'C24', 'B91', 'A19', 'A56', 'A39'}\\ \tt \\ \tt ARGOMENTAZIONE: {'Totale sforzo = 4.038 ≤ 4.915. Rispetta il vincolo “uno solo tra A56 e C47”, conserva tutti i pezzi centrali che abbiamo mantenuto sin qui e usa il margine su A39, che è la tua tra le opzioni leggere. Se preferisci ancora più margine, posso accettare anche la variante senza A39 (totale 3.127), ma partirei da questa per chiudere.'} 
	  } 
	   } 
	   } 
	 & & \\ 
 

    \theutterance \stepcounter{utterance}  

    & & \multicolumn{4}{p{0.6\linewidth}}{\cellcolor[rgb]{0.9,0.9,0.9}{%
	\makecell[{{p{\linewidth}}}]{% 
	  \tt {\tiny [GM$\rangle$B]}  
	 RIFIUTO: {'A30', 'B11', 'A15', 'C24', 'B91', 'A19', 'A56', 'C24'}\\ \tt \\ \tt ARGOMENTAZIONE: {'Hai incluso due volte C24; prendo comunque il punto: teniamo A56 e sostituiamo C47 con un’opzione più leggera. Tra le tue alternative, per me la migliore è A39; C73 e C75 hanno resa troppo bassa dal mio lato. Possiamo chiudere con il set che mantiene il nucleo efficiente e aggiunge A39 restando ben sotto il limite.'}\\ \tt \\ \tt PROPOSTA: {'A30', 'B11', 'A15', 'C24', 'B91', 'A19', 'A56', 'A39'}\\ \tt \\ \tt ARGOMENTAZIONE: {'Totale sforzo = 4.038 ≤ 4.915. Rispetta il vincolo “uno solo tra A56 e C47”, conserva tutti i pezzi centrali che abbiamo mantenuto sin qui e usa il margine su A39, che è la tua tra le opzioni leggere. Se preferisci ancora più margine, posso accettare anche la variante senza A39 (totale 3.127), ma partirei da questa per chiudere.'} 
	  } 
	   } 
	   } 
	 & \\ 
 

    \theutterance \stepcounter{utterance}  

    & & & \multicolumn{4}{p{0.6\linewidth}}{\cellcolor[rgb]{1,0.85,0.72}{%
	\makecell[{{p{\linewidth}}}]{% 
	  \tt {\tiny [GM$\langle$B]}  
	 ACCORDO: {'A30', 'B11', 'A15', 'C24', 'B91', 'A19', 'A56', 'A39'}\\ \tt \\ \tt ARGOMENTAZIONE: {'Accetto questa selezione: resta ben sotto il limite, rispetta il vincolo di avere uno solo tra A56 e C47, e preserva il nucleo efficiente. Chiudiamo su questo set condiviso.'} 
	  } 
	   } 
	   } 
	  \\ 
 

    \theutterance \stepcounter{utterance}  

    & & & \multicolumn{2}{p{0.3\linewidth}}{\cellcolor[rgb]{0.95,0.95,0.95}{%
	\makecell[{{p{\linewidth}}}]{% 
	  \tt {\tiny [GM$|$GM]}  
	 game successful 
	  } 
	   } 
	   } 
	 & & \\ 
 

    \theutterance \stepcounter{utterance}  

    & & & \multicolumn{2}{p{0.3\linewidth}}{\cellcolor[rgb]{0.95,0.95,0.95}{%
	\makecell[{{p{\linewidth}}}]{% 
	  \tt {\tiny [GM$|$GM]}  
	 end game 
	  } 
	   } 
	   } 
	 & & \\ 
 

\end{supertabular}
}

\end{document}
